\documentclass{article}
\usepackage{amsmath}
\begin{document}

\section{Основные уравнения UAF}

\subsection{Динамика когерентности уровня $l$}
\[
\frac{\partial \Gamma_l}{\partial \tau} = \eta_l \left( \sum_{k<l} \Phi_{k\to l} \mathcal{E}_{k,l} \right) (1-\alpha_l)^2 
- \lambda_l \Delta S_{env,l} \alpha_l \Gamma_l
\]

\subsection{Динамика интеграции}
\[
\frac{\partial \alpha_l}{\partial \tau} = \kappa_l \alpha_l (1-\alpha_l)(E_l - c_l) - \delta_l(1-\alpha_l)\Gamma_l^2 
+ \mu_l \sum_{k<l} \mathcal{E}_{k,l}(\Gamma_k - \Gamma_l)
\]

\subsection{Закон сохранения полного времени}
\[
t_{past} + t_{self} + t_{future} = \text{const}
\]

\subsection{Вероятность туннелирования в UAF-Q}
\[
P_{\text{tunnel}}(\tau) = \Gamma_0(\tau)^2 \exp\left(-\int_0^\tau \alpha_1(s)(1-\Gamma_0(s))\,ds\right)
\]

\subsection{Пример калибровки (muon-catalyzed fusion)}
\[
\text{Среднее число реакций на мюон} = \eta \cdot (1-\alpha)^2 - \lambda \alpha \approx 150
\]

\end{document}
